%% tags: [breadth first search]
%% source: problems
\begin{prob}
    The code below is a modified version of BFS shown in class. This
    modification has a single addition: a \mintinline{python}{print} statement. A comment
    in the code indicates where the \mintinline{python}{print} statement is located.

    \begin{minted}[autogobble]{python}
        def bfs(graph, source):
            """Start a BFS at `source`."""
            status = {node: 'undiscovered' for node in graph.nodes}
            status[source] = 'pending'
            pending = deque([source])

            while pending: 
                u = pending.popleft()
                for v in graph.neighbors(u):
                    print("here!") # <--- this is the print statement
                    if status[v] == 'undiscovered':
                        status[v] = 'pending'
                        pending.append(v)
                status[u] = 'visited'
    \end{minted}

    Suppose this modified BFS is run on the following graph, starting at node
    $a$.

    \includegraphics[width=.2\textwidth]{graph.png}

    How many times will \mintinline{python}{"here!"} be printed? Your answer should be an exact number.

    \inlineresponsebox[.75in]{10}

\end{prob}
